% Converted from: paper/sections/04_dynamic_model_and_hydrodynamic_closure.md
\section{Dynamic Model and Hydrodynamic Closure}

\subsection{Modeling Scope and Architecture of the Reduced-Order Description}

This section develops the dynamic model used to describe the horizontal-launch to vertical-stabilization transition defined in Section~2. The objective is not to reproduce every fluid--structure interaction detail directly, but to construct a reduced-order model that remains physically interpretable, computationally efficient, and compatible with the free-decay evidence governed in Section~3. To achieve this, the model combines three elements: a planar 3-DOF rigid-body/hydrostatic core, CFD-derived static hydrodynamic mappings for force and moment closure, and experimentally identified dynamic terms that cannot be inferred reliably from static CFD alone.

A key design principle is explicit equation governance. CFD-derived static coefficients are used to replace the corresponding static hydrodynamic force and moment mappings, while rotational damping remains an identification-oriented term and is not double counted through ad hoc additions. This governance rule is essential for preserving interpretability and for maintaining consistency between manuscript equations and the simulation implementation.

\subsection{Governing Equations for the Transition Phase}

Using the state interfaces $\nu=[u,w,q]^T$ and $\eta=[\theta]$, the simulation state is written as
\[
\mathbf{x}=[u,\,w,\,q,\,\theta]^T,
\]
with kinematic closure
\[
\dot{\theta}=q.
\]

For the transition phase of interest, the governing dynamics are expressed in a planar Fossen-style form with effective inertias $m_x$, $m_z$, and $I_y$:
\[
\begin{aligned}
 m_x\dot{u}-m_zwq &= X_{cfd}(\alpha) - (W-B)\sin\theta + T, \\
 m_z\dot{w}+m_xuq &= Z_{cfd}(\alpha) + (W-B)\cos\theta, \\
 I_y\dot{q} &= M_{cfd}(\alpha) + M_{damp}(q) + M_{bg}(\theta) + M_{cable}(\theta,q) + M_{thruster}, \\
 \dot{\theta} &= q.
\end{aligned}
\]

Here $W=m_{dry}g$ and $B=B_{mass}g$. The terms $T$ and $M_{thruster}$ denote optional actuation inputs. They are set to zero during free-decay identification and validation unless a specific actuation strategy is being studied in Section~7.

The buoyancy restoring moment is expressed as
\[
M_{bg}(\theta)=B\left(z_b\sin\theta + x_b\cos\theta\right),
\]
and the rotational damping term is modeled as
\[
M_{damp}(q)=-(d_q+d_{q,abs}|q|)q.
\]
When cable effects are enabled in the hardware configuration, the cable contribution is modeled as
\[
M_{cable}(\theta,q)=-K_{cable}(\theta-\theta_{eq})-C_{cable,q}q.
\]

These expressions match the runtime formulation implemented in the simulation module used for validation and parametric analysis.

\subsection{Hydrodynamic Closure and Angle-of-Attack Mapping}

Static hydrodynamic loads are closed through CFD-derived coefficient mappings as functions of the attack angle $\alpha$. The body-frame speed and dynamic pressure are defined as
\[
V=\sqrt{u^2+w^2}, \qquad Q=\tfrac{1}{2}\rho V^2,
\]
with a numerical safeguard for low-speed states to prevent unstable coefficient evaluation. The model then uses
\[
X_{cfd}=QA_{ref}C_X(\alpha), \qquad
Z_{cfd}=QA_{ref}C_Z(\alpha), \qquad
M_{cfd}=QA_{ref}L_{ref}C_m(\alpha).
\]

The attack angle is computed from the frozen sign convention,
\[
\alpha=\mathrm{atan2}(w,u),
\]
and is mapped to the available CFD lookup range at runtime. The coefficient convention is consistent across the manuscript and implementation: $C_X$ acts along $+x_b$, $C_Z$ acts along $+z_b$, and $C_m$ acts about $+y_b$, with positive $C_m$ corresponding to nose-up moment.

The CFD role in this paper is intentionally limited to static hydrodynamic closure and mechanism evidence. In particular, when the static pitch-moment mapping $C_m(\alpha)$ already reflects Munk-like static effects, no additional explicit Munk-moment term is added to the reduced-order pitch equation.

\subsection{Anisotropic Permeability-Corrected Added Mass and Effective Inertias}

A central modeling feature of this study is the anisotropic permeability correction used to represent directional coupling between rigid-body motion and inner-water inertia. Instead of treating inner-water effects as either fully decoupled or fully rigidly attached, the model introduces three channel-specific coupling parameters,
\[
\mu_x,\,\mu_z,\,\mu_{\theta} \in [0,1],
\]
which describe the fraction of inner-water inertia effectively coupled to surge, heave, and pitch motion, respectively. Values near zero indicate weak coupling (high effective internal permeability), whereas values near one indicate near-rigid coupling.

Using the project sign convention for added-mass derivatives, the total added terms are written as
\[
\begin{aligned}
X_{\dot{u},total} &= X_{\dot{u},outer} - \mu_x m_{water,inner}, \\
Z_{\dot{w},total} &= Z_{\dot{w},outer} - \mu_z m_{water,inner}, \\
M_{\dot{q},total} &= M_{\dot{q},outer} - \mu_{\theta} I_{water,inner}.
\end{aligned}
\]

The effective inertias used in the ODE denominators are then
\[
m_x=m_{dry}-X_{\dot{u},total}, \qquad m_z=m_{dry}-Z_{\dot{w},total}, \qquad I_y=I_{yy}-M_{\dot{q},total}.
\]

An important modeling choice, implemented consistently in the codebase, is that the translational denominators use $m_{dry}$ rather than $m_{wet}$. This convention is retained throughout the manuscript to avoid hidden discrepancies between equations, configuration files, and numerical simulations.

% TODO: Insert Fig. 4--Fig. 6 here (see paper/figures/FIGURE_REQUIREMENTS.md).

\subsection{CFD-Supported Coefficient Acquisition and Mechanism Interpretation}

The hydrodynamic coefficient mappings used above are obtained from CFD and treated as model-evidence inputs. Because the reduced-order simulation depends strongly on these mappings over a wide attitude range, the CFD reporting in this paper is organized to support both credibility and interpretability. The first requirement is numerical credibility through mesh-independence and residual-convergence summaries. The second is mechanism interpretability through coefficient trends, pressure-center migration, and representative field visualizations at selected attack angles.

The resulting CFD evidence package is therefore structured in four layers: numerical verification, coefficient mappings, pressure-center migration, and representative flow/pressure fields. This ordering is deliberate. It allows the reader to assess numerical trustworthiness before interpreting the physics encoded in the coefficient curves and moment trends used by the reduced-order model.

% TODO: Insert Fig. 7--Fig. 10 here (see paper/figures/FIGURE_REQUIREMENTS.md).

\subsection{Transition to Identification and Validation}

Section~4 defines the closed equation structure and the governed hydrodynamic closure logic. Parameter identification and uncertainty-oriented diagnostics are presented in Section~5. The resulting calibrated model package is then transferred without retuning to validation (Section~6) and to the bounded design guidance analyses (Section~7).
